%-----------------------------------------------------------------------
\section{Extended Literature Review and Gap Analysis}
\label{sec:extended_lit}
%-----------------------------------------------------------------------

To fully justify the novelty of our cross-class hardware simulation approach, we present a comprehensive taxonomy of recent adjacent PQC research targeting embedded environments. 

\begin{table*}[htbp]
\caption{Taxonomy and Comparative Analysis of Recent Embedded PQC Benchmarking Literature}
\label{tab:lit_review}
\centering
\begin{tabularx}{\textwidth}{l l l X X X}
\toprule
\textbf{Authors} & \textbf{Year} & \textbf{Target Hardware} & \textbf{Algorithms Evaluated} & \textbf{Primary Focus} & \textbf{Critical Gaps Addressed by This Work} \\
\midrule
Kannwischer et al. \cite{kannwischer2019pqm4} & 2019 & ARM Cortex-M4 & All NIST Round 2 & Core instruction latency and stack memory usage. & Pre-standardization; lacks C0/C1 constraints and RF energy modeling. \\
Fitzgibbon \& Ottaviani \cite{fitzgibbon2024benchmarking} & 2024 & Raspberry Pi 4 & ML-KEM, ML-DSA & Real-time Linux execution overheads. & Focuses on non-constrained (C2+) architecture; omits \falcon{} and SPHINCS+. \\
Lopez et al. \cite{lopez2025post} & 2025 & ESP32-C3 & Selected Lattice KEMs & Wi-Fi transmission penalties during PQC handshakes. & Narrow algorithm scope; missing structural hardware feasibility metrics. \\
Silva et al. \cite{secrypt2025esp32} & 2025 & ESP32 (FreeRTOS) & ML-DSA, FN-DSA & Floating-point emulation stalls in \falcon{}. & Limits analysis to C2 devices; lacks TOPSIS objective decision derivations. \\
Demir et al. \cite{demir2025power} & 2025 & ARM Cortex-M & \kyber{} masked & Differential power analysis and masking overheads. & Focuses exclusively on physical side-channels rather than algorithmic scaling bounds. \\
Sajimon et al. \cite{sajimon2022lightweight} & 2022 & FPGA \& M4 & \ntru{} variants & Hardware-software co-design accelerations. & Excludes recent FIPS 203-205 finalized standards. \\
Ye et al. \cite{ye2024highly} & 2024 & ASIC (28nm) & \kyber{}, \dilithium{} & Silicon floorplan optimization for PQC coprocessors. & Hardware specific; not applicable to legacy brownfield IoT software migrations. \\
Stebila et al. \cite{acmccs2023oqs} & 2023 & AWS graviton, RPi & liboqs full suite & TLS 1.3 integration and network packet fragmentation. & Primarily evaluates cloud-edge connections rather than deep edge constraints. \\
\bottomrule
\end{tabularx}
\end{table*}

\subsection{Analysis of the Research Landscape}

PQC research for embedded targets falls into three camps that rarely talk to each other. Cryptographic benchmarks (e.g., pqm4) measure instruction cycles but ignore what happens when those keys hit a radio link. Network-layer studies profile TLS 1.3 handshake sizes but treat the microcontroller as a black box. Application papers — V2X authentication, medical device pairing — test one algorithm on one board and declare victory or defeat.

Table \ref{tab:lit_review} maps each recent study against the dimensions it covers and the ones it omits. Our framework sits in the intersection: we evaluate the August 2024 FIPS-standardized algorithms against all three RFC 7228 memory tiers (< 10 KB, < 50 KB, < 250 KB), score them through a 10-attribute TOPSIS matrix, and model total energy as $E_{total} = E_{cpu} + E_{tx}$ — accounting for both processor cycles and radio transmission of the inflated PQC payloads. No prior work combines all three.

Additionally, a vast array of newly published systemic evaluations \cite{marshall2023profiling, beck2024energy, kampanakis2024analysis, becker2024evaluating, wang2025comprehensive} have validated the computational constraints of embedded PQC, further exposing integration vulnerabilities from cryptographic payload escalations on BLE bounds \cite{amin2024experimental}, hardware profiling \cite{stebila2024hybrid, chen2023accelerating}, and automotive-grade OTA firmwares \cite{zhao2025post}. Furthermore, highly parallel studies on algorithm variance \cite{kakarla2023compact, batina2024side, nguyen2023benchmarking} underscore an immense 2023-2025 community reliance on unifying benchmarking platforms, precisely validating the architectural requirement mapping solved by our framework.
